\subsection{Redshift Dependence of the Hubble Offset}
  \label{subsec:redshift_dependence}

  The environmental bias predicted by the framework is inherently scale dependent.

  At sufficiently low redshift, local measurements are sensitive to individual
  large-scale structures and directly probe the saturated regime.
  At increasing redshift, line-of-sight averaging over multiple environments
  progressively suppresses the bias.

  As a result, the inferred expansion rate is predicted to converge smoothly
  toward $H_{\mathrm{global}}$ with increasing redshift.

  This behavior provides a clear observational signature.
  Measurements of the Hubble constant performed over different redshift ranges
  should exhibit a systematic trend, with the largest deviations appearing at the
  lowest redshifts, followed by a gradual convergence at higher
  redshift~\cite{Verde2019Review}.

  An important quantitative consequence of the framework is that the expected
  deviation remains at the percent level.
  Once the saturation scale $a_\star$ is fixed by galactic dynamics, the fractional
  offset in the inferred expansion rate is predicted to peak at intermediate
  redshift and to remain negligible at both very low and sufficiently high redshift.

  The peak is expected when the typical line-of-sight length becomes comparable
  to the characteristic scale of large voids, before full environmental averaging
  sets in.

  Such deviations fall within the sensitivity range of forthcoming large-scale
  surveys, including DESI and \textit{Euclid}.
  Observed deviations significantly larger or smaller than this level would
  directly challenge the proposed mechanism.
